\begin{frame}
    \frametitle{Math}
    \framesubtitle{More functions}
    % [Speakers Notes]
    % But wait, this looks weird!?
    % This is because we wrote this in inline-math mode. 
    % If we write this in the displaymath environment, 
    % the limits will appear below and above the symbols, as we would expect.
    % If we desire this behavior in inline math as well, 
    % there's a fix for this: The \limits command.
    \note{%
    \begin{itemize}
        \item This looks weird\ldots
        \item Because of inline-math mode (tries to minimize vertical size)
        \item We need displaymath behavior 
        \item Solution: limits command
    \end{itemize}}
    \begin{center}
        \Large%
        But wait, this looks weird!?

        \vfill
        $\lim_{n \to \infty} \frac{1}{n}
        \qquad
        \sum_{i = 0}^{n} a_n
        \qquad
        \prod_{j = 0}^{m}b_m
        \qquad
        \int_{a}^{b} x d\,x$

        \vfill
        \only<2->{There's a fix for this!}
    \end{center}
\end{frame}

\begin{frame}[fragile]
    \frametitle{Math}
    \framesubtitle{\texttt{\bslash limits} command}
    \begin{center}
        \Large%
        The \lstinline|\limits| command.
    \end{center}
    \begin{columns}
        \begin{column}{.5\textwidth}
            \begin{exlisting}{Example}
A different inline sum
$\sum\limits_{i = 0}^{n} a_n$
is easy.    
            \end{exlisting}
        \end{column}
        \begin{column}{.48\textwidth}
            \centering
            \vspace{1em}
            
            A different inline sum
            $\sum\limits_{i = 0}^{n} a_n$
            is easy.
        \end{column}
    \end{columns}
\end{frame}

\begin{frame}
    \frametitle{Math}
    \framesubtitle{Custom functions}
        {\Large\lstinline|\\operatorname\{definition\}|}
        
        \begin{itemize}
            \item<+-> Typesets contents as a function / an operator
            \item<+-> Handles spacing correctly
            \item<+-> Should only be used if function is only used once
        \end{itemize}
        \pause
        
        \vspace{1.5em}
        {\Large\lstinline|\\DeclareMathOperator\{command\}\{definition\}|}
        
        \begin{itemize}
            \item Defines new math function
            \item Can \textbf{only} be used in the \emph{preamble}!
        \end{itemize}
\end{frame}

\begin{frame}[fragile]
    \frametitle{Math}
    \framesubtitle{Custom functions}
    \begin{columns}
        \begin{column}{.6\textwidth}
            \begin{exlisting}{Example}
\DeclareMathOperator{\arcsinh}{arcsinh}
$\operatorname{arccosh}x + \arcsinh y$
            \end{exlisting}
        \end{column}
        \begin{column}{.38\textwidth}
            \centering
            \vspace{1em}

            $\operatorname{arccosh}x + \arcsinh y$
        \end{column}
    \end{columns}
\end{frame}

\begin{frame}
    \frametitle{Exercise}
    \begin{columns}
        \begin{column}{.48\textwidth}
            \begin{align*}
                x_{t+1} 
                    &= kx_t (1 - x_t) \\[1.2em]
                H 
                    &= - \sum_{i = 0}^{N} p(x_i) \log p(x_i)
            \end{align*}
            \begin{center}
                Advanced:
            \end{center}
            \begin{align*}
                \mathcal{L}^{-1}\{F\} 
                    &=  \lim_{T \to \infty}
                        \frac{1}{2\pi i}
                        \int\limits_{\gamma - iT}^{\gamma + iT}
                        e^{st} F(s) \,\mathrm{d}s
            \end{align*}
        \end{column}
        \begin{column}{.48\textwidth}
            \begin{block}{\blockstrut Resources}
                \centering
                List of \LaTeX{} Mathematical Symbols

                \medskip
                \qrcode[%
                    level=L, 
                    version=0,
                    height=2cm]{https://oeis.org/wiki/List_of_LaTeX_mathematical_symbols}
                
                
                \bigskip
                Overleaf Documentation
                
                \medskip
                \qrcode[%
                    level=L, 
                    version=0,
                    height=2cm]{https://www.overleaf.com/learn}                
            \end{block}
        \end{column}
    \end{columns}
\end{frame}